\scanpage{18}
\uline{Lemme 1} Un morphisme \deleted{radiciel surjectif}\inserted{de préschémas}
$f:T\mathbin{\text{---}}S$ qui est un homéomorphisme \unsure{sur|uni}
\deleted{\unsure{universel}} est affine.
\margin{\figure{\begin{tikzpicture}\draw (0,0)--(0,-1.1);\end{tikzpicture}}}

\deleted{En effet, \unsure{c'est} un homéomorphisme, il est
quasi-compact, Étant radiciel, il est séparé, donc
\deleted{$f_*(\mathcal O_T)=A$ est quasi-cohérent}
\inserted{pour tout $F$ sur $T$, $f_*(F)$ est} quasi-cohérent.
Supposons $S$ \deleted{affine, \unsure{prenons, pour voir les} $F$ qu. coh. sur $T$}
On peut supposer $S$ affine, donc quasi-compact, donc $T$ est quasi-compact.
On a donc $T=\bigcup_i T_i$, réunion \unsure{finie}
d'ouverts affines; \deleted{et $T_i$ est} comme
$f(T_i)$ est ouvert de $S$, il est réunion d'ouverts affines $S_{ij}$,
\deleted{\unsure{dont chacun}}
\inserted{et $T_i=\bigcup_j T_{ij}$, où $T_{ij}=f^{-1}(S_{ij})=f_i^{-1}(S_{ij})$}
est affine.
\deleted{$T=\bigcup_{ij}f^{-1}(S_{ij})$}
comme $T$ est quasi-cpt, $T$ est réunion finie}

On peut supposer $S$ \uline{affine}.
Soit $s\in S$, alors $f^{-1}(s)$ est réduit à un pt,
donc admet un voisinage affine $U$.
Comme $f(U)$ est un voisinage de $s$, il contient un
\uline{voisinage} affine $V$ de $s$. Alors on a
\[
f^{-1}(V)=f'^{-1}(V)
\]
où $f':U\mathbin{\text{---}}S$ est induit par $f$.
Comme $f'$ est un morphisme de schémas affines, et
$V$ est affine, $f'^{-1}(V)$ l'est, O.K.
\mspage{17}
